\section{哈密顿扰动\personCauchy{}--\personRiemann{}方程}
\label{sec:hamiltonian-perturbed-cauchy-riemann-equation}

为了对 $\mathcal A_H$ 作\personMorse{}理论，需要在
$\widetilde{\mathcal L}_0(M)$ 上给出度量。为此先在 $\mathcal L_0(M)$ 上定义度量，
再将其拉回 $\widetilde{\mathcal L}_0(M)$。注意，$M$ 上任何\personRiemann{}度量的
$S^1$-族 $\{g_t\}_{t\in S^1}$ 都在 $\mathcal L(M)$ 上诱导 $L^2$ 型度量
\begin{equation}
 \langle\!\langle\xi_1,\xi_2\rangle\!\rangle
 =\int_0^1g_t(\xi_1(t),\xi_2(t))\,dt
 \label{eq:18-loop-space-l2-metric}
\end{equation}
（$\xi_1,\xi_2\in T_\gamma\mathcal L(M)$）。在辛流形 $(M,\omega)$ 上，特别考虑由
\stfhIndex{I034}{Daikaileduliang}{殆凯勒度量（almost-Kähler metric）}
与辛形式 $\omega$ 相容的殆复结构一参数族 $J$ 所诱导的\emph{殆\personKahler{}度量}族
\[
 g_J=\omega(\,\cdot,J\,\cdot)
\]
及相伴范数 $|\cdot|_J$。给定 $J=\{J_t\}_{t\in S^1}$，它在 $\mathcal L(M)$ 上诱导
相伴的 $L^2$ 度量 $\langle\!\langle\cdot,\cdot\rangle\!\rangle_J$，也可写成
\begin{equation}
 \langle\!\langle\xi_1,\xi_2\rangle\!\rangle_J
 =\int_0^1\omega(\xi_1(t),J_t\xi_2(t))\,dt
 \label{eq:18-loop-space-compatible-l2-metric}
\end{equation}
（$\xi_1,\xi_2\in T_\gamma\mathcal L(M)$）。除非另作说明，从现在起 $J$ 总表示相容
殆复结构的 \emph{$S^1$-族}，并记
\stfhIndex{I010}{00-jomega}{$\mathfrak j_\omega$}
\[
 \mathfrak j_\omega:=C^\infty(S^1,\mathcal J_\omega).
\]
若以 $\operatorname{grad}_J\mathcal A_H$ 表示相伴的 $L^2$ 梯度向量场，则由
\eqref{eq:18-global-action-primitive} 与
\eqref{eq:18-loop-space-compatible-l2-metric} 直接计算得到
\begin{equation}
 \operatorname{grad}_J\mathcal A_H([\gamma,w])(t)
 =J_t\bigl(\dot\gamma(t)-X_H(t,\gamma(t))\bigr),
 \label{eq:18-action-l2-gradient}
\end{equation}
以后简写为 $J(\dot\gamma-X_H(\gamma))$。

\stfhSubIndexRangeEnd{I065}{Kangli-Zengdeerzhibiao}{康利--曾德尔指标（Conley--Zehnder index）}{Zzwdezhibiao}{$[z,w]$ 的指标（of $[z,w]$）}
\begin{exercise}
\label{ex:18-prove-action-gradient}
证明这个公式。
\end{exercise}

由此可知，该梯度可以投影到 $\mathcal L_0(M)$。因此，把道路
$u:\mathbb R\to\widetilde{\mathcal L}_0(M)$ 的\emph{负}梯度流方程投影到
$\mathcal L_0(M)$ 后，若把 $u$ 看作映射 $u:\mathbb R\times S^1\to M$，方程写成
\begin{equation}
 \frac{\partial u}{\partial\tau}
 +J\left(\frac{\partial u}{\partial t}-X_H(u)\right)=0.
 \label{eq:18-floer-perturbed-cr-equation}
\end{equation}
称它为与二元组 $(H,J)$ 相伴的\emph{\personFloer{} \termPerturbedCauchyRiemannEquation{}}，或简称
扰动\personCauchy{}--\personRiemann{}方程。

\personFloer{}理论在很大程度上依赖于研究形如
\eqref{eq:18-floer-perturbed-cr-equation} 的扰动\personCauchy{}--\personRiemann{}方程的\emph{有限能量}解
$u:\mathbb R\times S^1\to M$ 的模空间。相应的离壳能量泛函定义如下。

\begin{definition}[能量]
\label{def:18-floer-energy}
给定光滑映射 $u:\mathbb R\times S^1\to M$，定义 $u$ 的能量
$E_{(H,J)}(u)$ 为
\[
 E_{(H,J)}(u)=\frac12\int_{\mathbb R\times S^1}
 \left(\left|\frac{\partial u}{\partial\tau}\right|_{J_t}^2
 +\left|\frac{\partial u}{\partial t}-X_H(u)\right|_{J_t}^2\right)
 dt\,d\tau.
\]
\end{definition}

下面的命题是引理 12.2.4 的闭弦版本，证明与此前相同，故略去。

\begin{proposition}
\label{prop:18-finite-energy-asymptotic-subsequence}
设 $(M,\omega)$ 是任意辛流形，不必紧，$H:S^1\times M\to\mathbb R$ 是任意
哈密顿量。设 $u:\mathbb R\times S^1\to M$ 是
\eqref{eq:18-floer-perturbed-cr-equation} 的有限能量解，且其像在 $\operatorname{int}M$
中有紧闭包。则存在序列 $\tau_k\to\infty$（分别地，$\tau_k\to-\infty$），使环路
$z_k:=u(\tau_k)=u(\tau_k,\cdot)$ 在 $C^\infty$ 意义下收敛到\personHamilton{}方程
$\dot x=X_H(x)$ 的一个 1-周期解 $z:S^1\to M$。
\end{proposition}

对不必非退化的一般 $H$，记
\[
 \mathcal M(H,J)=\mathcal M(H,J;\omega)
\]
为~\eqref{eq:18-floer-perturbed-cr-equation} 的有限能量解的集合。

对~\eqref{eq:18-floer-perturbed-cr-equation} 的\emph{非自治}版本可以作类似讨论，现说明如下。
首先回忆
\[
 \mathcal H_0(M)=\{H:S^1\times M\to\mathbb R\mid H\text{ 均值归一化}\}.
\]
考虑 $\mathbb R$-族
\[
 \mathcal H_R:\mathbb R\to\mathcal H_0,\quad\tau\mapsto H(\tau),
 \qquad
 j_R:\mathbb R\to\mathfrak j_\omega,\quad\tau\mapsto J(\tau),
\]
它们渐近常值，即当 $|\tau|>R$（其中 $R$ 是充分大的常数）时，
\[
 H(\tau)=H^{\pm\infty},\qquad J(\tau)=j^{\pm\infty},
\]
其中 $H^{\pm\infty}\in\mathcal H_0$、$j^{\pm\infty}\in\mathfrak j_\omega$。
\footnote{译者注：源书印 $\mathbb J_\mathbb R\colon\mathbb R\to\mathcal J_\omega$、
$j^{\pm\infty},j\in\mathcal J_\omega$、$J_\alpha,J_\beta\in\mathcal J_\omega$；
译文统一作 $j_\mathbb R$、$j^{\pm\infty}$ 等，且陪域归 $\mathfrak j_\omega$——因
式 (18.5.42) 与 (18.5.44) 中的 $J(\tau)$ 须为 $S^1$ 族，并须与第 19 章
$(\mathcal H^\rho,\mathfrak j^\rho)$ 相衔接。}
与任意
这样的二元组相伴有~\eqref{eq:18-floer-perturbed-cr-equation} 的如下\emph{非自治}版本：
\begin{equation}
 \frac{\partial u}{\partial\tau}
 +J(\tau)\left(\frac{\partial u}{\partial t}-X_{H(\tau)}(u)\right)=0.
 \label{eq:18-nonautonomous-floer-equation}
\end{equation}
相伴能量泛函为
\[
 E_{(\mathcal H_R,j_R)}(u)=\frac12\int_{-\infty}^{\infty}\int_0^1
 \left(\left|\frac{\partial u}{\partial\tau}\right|_{J(\tau)}^2
 +\left|\frac{\partial u}{\partial t}-X_{H(\tau)}(u)\right|_{J(\tau)}^2\right)
 dt\,d\tau.
\]
记
\[
 \mathcal M(\mathcal H_R,j_R)=\mathcal M(\mathcal H_R,j_R;\omega)
\]
为~\eqref{eq:18-nonautonomous-floer-equation} 的有限能量解的集合。

下面是命题~\ref{prop:18-finite-energy-asymptotic-subsequence} 的类似结果。由于
$(\mathcal H,j)$ 渐近常值，其证明本质上相同。

\begin{proposition}
\label{prop:18-nonautonomous-asymptotic-subsequence}
设 $\mathcal H_R,j_R$ 如上，且 $u:\mathbb R\times S^1\to M$ 是
\eqref{eq:18-nonautonomous-floer-equation} 的有限能量解。则存在序列
$\tau_k\to\infty$，使环路 $z_k:=u(\tau_k)=u(\tau_k,\cdot)$ 在 $C^\infty$ 意义下
收敛到\personHamilton{}方程 $\dot x=X_{H^{+\infty}}(x)$ 的 1-周期解
$z:S^1\to M$。在 $-\infty$ 端也有相同结论，此时极限满足
$\dot x=X_{H^{-\infty}}(x)$。
\end{proposition}

这类渐近常值族的一种典型来源，是把 $[0,1]$ 上给定的光滑一参数族\emph{拉长}。

\begin{definition}
\label{def:18-boundary-flat-map}
设 $T$ 是任意拓扑空间。若连续映射 $f:[0,1]\to T$ 在边界
\stfhIndex{I044}{Bianjiepingtan}{边界平坦（boundary flat）}
$\partial[0,1]=\{0,1\}$ 的邻域内为常值，则称 $f$ \emph{边界平坦}。
\end{definition}

设 $\mathcal H:[0,1]\to\mathcal H_0(M)$ 是连接两个哈密顿量
$H_\alpha,H_\beta\in\mathcal H_0(M)$ 的同伦，
$j:[0,1]\to\mathfrak j_\omega$ 连接 $J_\alpha,J_\beta\in\mathfrak j_\omega$。记
\[
 \mathcal P(\mathfrak j_\omega):=C^\infty([0,1],\mathfrak j_\omega),
 \qquad
 \mathcal P(\mathcal H_0(M)):=C^\infty([0,1],\mathcal H_0(M)).
\]
定义一类函数 $\rho:\mathbb R\to[0,1]$：
\begin{equation}
 \rho(\tau)=\begin{cases}
 0,&\tau\leq0,\\
 1,&\tau\geq1,
 \end{cases}
 \qquad \rho'(\tau)\geq0.
 \label{eq:18-elongation-function}
\end{equation}
\stfhIndex{I072}{Lachanghanshu}{拉长函数（elongation function）}
称 $\rho$ 为\termElongationFunction{}。

每个这样的二元组 $(\mathcal H,j)$ 连同拉长函数 $\rho$，定义一对渐近常值
$\mathbb R$-族 $(\mathcal H^\rho,j^\rho)$：
\[
 \mathcal H_R=\mathcal H^\rho,\qquad j_R=j^\rho,
\]
其中 $\mathcal H^\rho=\{H^\rho(\tau)\}_{\tau\in\mathbb R}$ 是重新参数化的同伦
\[
 \tau\longmapsto H^\rho(\tau,t,x)=H(\rho(\tau),t,x).
\]
称 $\mathcal H^\rho$ 为 $\mathcal H$ 的 \emph{$\rho$-拉长}，或 $\mathcal H$ 的
\emph{$\rho$-拉长同伦}；对 $j$ 作相同定义。因此，这样的三元组 $(\mathcal H,j;\rho)$ 给出
\emph{非自治}方程
\begin{equation}
 \frac{\partial u}{\partial\tau}
 +J^\rho\left(\frac{\partial u}{\partial t}-X_{H^\rho}(u)\right)=0.
 \label{eq:18-elongated-floer-equation}
\end{equation}
记
\[
 \mathcal M(\mathcal H,j;\rho)
\]
为~\eqref{eq:18-elongated-floer-equation} 的有限能量解集合。
